\centerline{\bf \Huge Клеточки)0))))))))))00))0}

\begin{flushright}
        \scriptsize{}
\end{flushright}


\problem{1}
Можно ли так расставить фишки в клетках доски $8 \times 8$, чтобы в каждых двух столбцах количество фишек было одинаковым, а в каждых двух строках - различным?

\problem{2}
Каждую грань кубика разбили на четыре одинаковых квадрата, а затем раскрасили эти квадраты в несколько цветов так, что квадраты, имеющие общую сторону, оказались окрашенными в разные цвета. Какое наибольшее количество квадратов цвета маренго могло получится?

\problem{3}
Под одной из клеток доски $8 \times 8$ зарыт клад. Под каждой из остальных зарыта табличка, в которой указано, за какое наименьшее число шагов можно добраться из этой клетки до клада, если одним шагом можно перейти из клетки в соседнюю по стороне. Какое наименьшее число клеток надо перекопать, чтобы наверняка достать клад?

\problem{4}
На шахматной доске стоят несколько ферзей. Известно, что каждый из них бьёт ровно $N$ других. Какие значения может принимать $N$ ?

\problem{5}
На шахматной доске стоит 33 коня. Докажите, что какой-то конь бьёт хотя бы двух других.